
\documentclass{article}
\usepackage{colortbl}
\usepackage{makecell}
\usepackage{multirow}
\usepackage{supertabular}

\begin{document}

\newcounter{utterance}

\centering \large Interaction Transcript for game `function\_detective', experiment `string\_passive\_labeled\_pairs', episode 129 with mistral{-}large{-}3{-}azure.
\vspace{24pt}

{ \footnotesize  \setcounter{utterance}{1}
\setlength{\tabcolsep}{0pt}
\begin{supertabular}{c@{$\;$}|p{.15\linewidth}@{}p{.15\linewidth}p{.15\linewidth}p{.15\linewidth}p{.15\linewidth}p{.15\linewidth}}
    \# & \multicolumn{2}{c}{Player} && \multicolumn{2}{c}{Game Master} \\
    \hline

    \theutterance \stepcounter{utterance}  
    & & & \multicolumn{4}{p{0.6\linewidth}}{
        \cellcolor[rgb]{0.9,0.9,0.9}{
            \makecell[{{p{\linewidth}}}]{
                \texttt{\tiny{[P1$\langle$GM]}}
                \texttt{I have a hidden Python function for you to figure out.} \\
\\ 
\texttt{Your job is to infer what it does and then submit a correct Python implementation.} \\
\\ 
\texttt{Function signature:} \\
\texttt{(text: str) {-}> str} \\
\\ 
\texttt{You have 15 turns total.} \\
\\ 
\texttt{I reveal candidate input{-}output pairs one at a time and tell you whether each one is correct.} \\
\\ 
\texttt{You must make your best SOLVE attempt no later than the second{-}to{-}last turn.} \\
\texttt{Do not wait until the very end.} \\
\\ 
\texttt{You can use these actions:} \\
\texttt{{-} JSON format described at the end of the prompt} \\
\texttt{{-} SOLVE: ```python ... ```} \\
\\ 
\texttt{Do not add import statements in your guess.} \\
\texttt{Anything you might need is already available, so just write the function body in hypothesis and return the answer.} \\
\\ 
\texttt{Preloaded examples:} \\
\texttt{1. (('a', ''), True)} \\
\\ 
\texttt{Scoring:} \\
\texttt{{-} correctness is required for any non{-}zero score} \\
\texttt{{-} in interactive modes, solving with fewer information requests is better} \\
\\ 
\texttt{Input formatting reminders:} \\
\texttt{{-} strings must be quoted} \\
\texttt{{-} booleans must be True/False} \\
\texttt{{-} use commas between multiple values} \\
\texttt{{-} do not write any extra action tags in the same message} \\
\\ 
\texttt{Format example:} \\
\texttt{You write:} \\
\texttt{\{'current\_mode': <probing/confirming/don't know; choose whichever best describes your current state in the task>, 'input\_rationale': <brief description of why you chose this hypothesis>, 'hypothesis': <either the exact string "unknown" if you have no guess yet OR your current best guess as runnable Python matching the given signature, with a one{-}line docstring stating what it does. Never write placeholder code, if you cannot commit to a guess, use "unknown". Match the given signature.} \\
\texttt{e.g.} \\
\texttt{def f(x):} \\
\texttt{	"""Returns the absolute value of x."""} \\
\texttt{	return x if x >= 0 else {-}x} \\
\\ 
\texttt{Use "unknown" if you have no guess.>, 'confidence\_score': <score between 1{-}100. 100 means highest confidence in your hypothesis>\}} \\
\\ 
\texttt{Output only one valid json per turn.} \\
            }
        }
    }
     \\ \\

    \theutterance \stepcounter{utterance}  
    & \multicolumn{4}{p{0.6\linewidth}}{
        \cellcolor[rgb]{0.561,0.737,0.561}{
            \makecell[{{p{\linewidth}}}]{
                \texttt{\tiny{[P1$\rangle$GM]}}
                \texttt{```json} \\
\texttt{\{} \\
\texttt{  "current\_mode": "probing",} \\
\texttt{  "input\_rationale": "The first example shows that when the input is 'a', the output is an empty string. This suggests the function might be removing or filtering certain characters from the input string. A common operation that results in an empty string for a single character is removing all vowels, but 'a' is a vowel, so this might not fit. Alternatively, it could be removing all non{-}vowel characters, but 'a' is a vowel, so the output would be 'a', not empty. Another possibility is removing all lowercase letters, but 'a' is lowercase, so the output would be empty, which matches. To test this, I'll hypothesize that the function removes all lowercase letters from the input string.",} \\
\texttt{  "hypothesis": "def f(text: str) {-}> str:\textbackslash{}n    \textbackslash{}"\textbackslash{}"\textbackslash{}"Returns the input string with all lowercase letters removed.\textbackslash{}"\textbackslash{}"\textbackslash{}"\textbackslash{}n    return ''.join({[}c for c in text if not c.islower(){]})",} \\
\texttt{  "confidence\_score": 70} \\
\texttt{\}} \\
\texttt{```} \\
            }
        }
    }
    & & \\ \\

    \theutterance \stepcounter{utterance}  
    & & & \multicolumn{4}{p{0.6\linewidth}}{
        \cellcolor[rgb]{0.9,0.9,0.9}{
            \makecell[{{p{\linewidth}}}]{
                \texttt{\tiny{[P1$\langle$GM]}}
                \texttt{OUTPUT: (('aA', '\_alt'), False)} \\
            }
        }
    }
     \\ \\

    \theutterance \stepcounter{utterance}  
    & \multicolumn{4}{p{0.6\linewidth}}{
        \cellcolor[rgb]{0.561,0.737,0.561}{
            \makecell[{{p{\linewidth}}}]{
                \texttt{\tiny{[P1$\rangle$GM]}}
                \texttt{```json} \\
\texttt{\{} \\
\texttt{  "current\_mode": "probing",} \\
\texttt{  "input\_rationale": "The previous hypothesis (removing all lowercase letters) fails for the input 'aA' because it would return 'A', but the actual output is '\_alt'. This suggests the function is not simply removing lowercase letters. The output '\_alt' hints at some kind of transformation or replacement, possibly related to case or special encoding. Another possibility is that the function is converting the string into a 'slug' or URL{-}safe format, replacing certain characters with underscores or other symbols. Alternatively, it might be a cipher or a specific encoding scheme. I'll now hypothesize that the function converts the input into a 'snake\_case' or 'slug' format, replacing non{-}alphanumeric characters and spaces with underscores, and converting uppercase letters to lowercase with a preceding underscore (e.g., 'A' becomes '\_a').",} \\
\texttt{  "hypothesis": "def f(text: str) {-}> str:\textbackslash{}n    \textbackslash{}"\textbackslash{}"\textbackslash{}"Converts the input string into a snake\_case/slug format, replacing non{-}alphanumeric characters with underscores and converting uppercase letters to lowercase with a preceding underscore.\textbackslash{}"\textbackslash{}"\textbackslash{}"\textbackslash{}n    result = {[}{]}\textbackslash{}n    for c in text:\textbackslash{}n        if c.isupper():\textbackslash{}n            result.append('\_' + c.lower())\textbackslash{}n        elif c.isalnum():\textbackslash{}n            result.append(c)\textbackslash{}n        else:\textbackslash{}n            result.append('\_} \\
            }
        }
    }
    & & \\ \\

    \theutterance \stepcounter{utterance}  
    & & & \multicolumn{4}{p{0.6\linewidth}}{
        \cellcolor[rgb]{0.9,0.9,0.9}{
            \makecell[{{p{\linewidth}}}]{
                \texttt{\tiny{[P1$\langle$GM]}}
                \texttt{Game Violation: ParseError: Response is not valid JSON.} \\
\texttt{Output ONLY:} \\
\texttt{One JSON with confidence\_score, current\_mode, input\_rationale, and hypothesis keys ...} \\
\texttt{OR} \\
\texttt{SOLVE: ```python ... ```} \\
            }
        }
    }
     \\ \\

    \theutterance \stepcounter{utterance}  
    & \multicolumn{4}{p{0.6\linewidth}}{
        \cellcolor[rgb]{0.561,0.737,0.561}{
            \makecell[{{p{\linewidth}}}]{
                \texttt{\tiny{[P1$\rangle$GM]}}
                \texttt{```json} \\
\texttt{\{} \\
\texttt{  "current\_mode": "probing",} \\
\texttt{  "input\_rationale": "The previous output '\_alt' for input 'aA' suggests the function is not a simple case transformation or slug conversion. The '\_alt' suffix hints at an alternative representation, possibly related to keyboard layouts (e.g., QWERTY to DVORAK) or some other substitution cipher. Given that 'a' alone returns an empty string, but 'aA' returns '\_alt', it might be checking for specific patterns or sequences. Another possibility is that the function is encoding modifier keys (e.g., 'Alt' for 'A'). I'll hypothesize that the function replaces uppercase letters with their corresponding modifier key representation (e.g., 'A' becomes '\_alt') and removes lowercase letters unless they are part of a sequence.",} \\
\texttt{  "hypothesis": "def f(text: str) {-}> str:\textbackslash{}n    \textbackslash{}"\textbackslash{}"\textbackslash{}"Replaces uppercase letters with their modifier key representation (e.g., 'A' {-}> '\_alt') and removes lowercase letters.\textbackslash{}"\textbackslash{}"\textbackslash{}"\textbackslash{}n    modifier\_map = \{'A': '\_alt', 'B': '\_ctrl', 'C': '\_ctrl', 'D': '\_alt', 'E': '\_alt', 'F': '\_shift', 'G': '\_shift', 'H': '\_shift', 'I': '\_alt', 'J': '\_shift', 'K': '\_shift', 'L': '\_shift', 'M': '\_} \\
            }
        }
    }
    & & \\ \\

    \theutterance \stepcounter{utterance}  
    & & & \multicolumn{4}{p{0.6\linewidth}}{
        \cellcolor[rgb]{0.9,0.9,0.9}{
            \makecell[{{p{\linewidth}}}]{
                \texttt{\tiny{[P1$\langle$GM]}}
                \texttt{Game Violation: ParseError: Response is not valid JSON.} \\
\texttt{Output ONLY:} \\
\texttt{One JSON with confidence\_score, current\_mode, input\_rationale, and hypothesis keys ...} \\
\texttt{OR} \\
\texttt{SOLVE: ```python ... ```} \\
            }
        }
    }
     \\ \\

    \theutterance \stepcounter{utterance}  
    & \multicolumn{4}{p{0.6\linewidth}}{
        \cellcolor[rgb]{0.561,0.737,0.561}{
            \makecell[{{p{\linewidth}}}]{
                \texttt{\tiny{[P1$\rangle$GM]}}
                \texttt{```json} \\
\texttt{\{} \\
\texttt{  "current\_mode": "probing",} \\
\texttt{  "input\_rationale": "Given the input 'aA' produces '\_alt', this strongly suggests the function is replacing uppercase 'A' with '\_alt' while removing lowercase 'a'. This resembles a keyboard shortcut encoding where uppercase letters are replaced by their modifier keys (e.g., 'Alt+A' becomes '\_alt'). The function likely removes lowercase letters and replaces uppercase letters with their modifier key representation. I'll refine the hypothesis to focus only on 'A' mapping to '\_alt' for now, as other mappings are unclear.",} \\
\texttt{  "hypothesis": "def f(text: str) {-}> str:\textbackslash{}n    \textbackslash{}"\textbackslash{}"\textbackslash{}"Replaces uppercase 'A' with '\_alt' and removes all other characters.\textbackslash{}"\textbackslash{}"\textbackslash{}"\textbackslash{}n    return '\_alt' if 'A' in text else ''",} \\
\texttt{  "confidence\_score": 80} \\
\texttt{\}} \\
\texttt{```} \\
            }
        }
    }
    & & \\ \\

    \theutterance \stepcounter{utterance}  
    & & & \multicolumn{4}{p{0.6\linewidth}}{
        \cellcolor[rgb]{0.9,0.9,0.9}{
            \makecell[{{p{\linewidth}}}]{
                \texttt{\tiny{[P1$\langle$GM]}}
                \texttt{OUTPUT: (('   ', '   '), True)} \\
            }
        }
    }
     \\ \\

    \theutterance \stepcounter{utterance}  
    & \multicolumn{4}{p{0.6\linewidth}}{
        \cellcolor[rgb]{0.561,0.737,0.561}{
            \makecell[{{p{\linewidth}}}]{
                \texttt{\tiny{[P1$\rangle$GM]}}
                \texttt{```json} \\
\texttt{\{} \\
\texttt{  "current\_mode": "probing",} \\
\texttt{  "input\_rationale": "The new example shows that whitespace is preserved exactly as{-}is. This contradicts my previous hypothesis (which would have removed spaces). The function must have special handling for whitespace. Combining this with earlier observations: (1) 'a' → '', (2) 'aA' → '\_alt', (3) '   ' → '   ', I now hypothesize that the function replaces uppercase 'A' with '\_alt', removes lowercase letters (except in special cases), and preserves whitespace. This suggests a keyboard shortcut encoder that only transforms specific uppercase letters while leaving other characters (including whitespace) untouched.",} \\
\texttt{  "hypothesis": "def f(text: str) {-}> str:\textbackslash{}n    \textbackslash{}"\textbackslash{}"\textbackslash{}"Replaces uppercase 'A' with '\_alt', removes lowercase letters, and preserves all other characters (including whitespace).\textbackslash{}"\textbackslash{}"\textbackslash{}"\textbackslash{}n    result = {[}{]}\textbackslash{}n    for c in text:\textbackslash{}n        if c == 'A':\textbackslash{}n            result.append('\_alt')\textbackslash{}n        elif c.islower():\textbackslash{}n            continue\textbackslash{}n        else:\textbackslash{}n            result.append(c)\textbackslash{}n    return ''.join(result)",} \\
\texttt{  "confidence\_score": 85} \\
\texttt{\}} \\
\texttt{```} \\
            }
        }
    }
    & & \\ \\

    \theutterance \stepcounter{utterance}  
    & & & \multicolumn{4}{p{0.6\linewidth}}{
        \cellcolor[rgb]{0.9,0.9,0.9}{
            \makecell[{{p{\linewidth}}}]{
                \texttt{\tiny{[P1$\langle$GM]}}
                \texttt{OUTPUT: (('a{-}b', '\_alt'), False)} \\
            }
        }
    }
     \\ \\

    \theutterance \stepcounter{utterance}  
    & \multicolumn{4}{p{0.6\linewidth}}{
        \cellcolor[rgb]{0.561,0.737,0.561}{
            \makecell[{{p{\linewidth}}}]{
                \texttt{\tiny{[P1$\rangle$GM]}}
                \texttt{```json} \\
\texttt{\{} \\
\texttt{  "current\_mode": "probing",} \\
\texttt{  "input\_rationale": "The input 'a{-}b' produces '\_alt', which contradicts my previous hypothesis (which would have removed 'a' and '{-}' while keeping 'b', resulting in 'b'). This suggests the function is not just removing lowercase letters and replacing 'A' with '\_alt'. The output '\_alt' for 'a{-}b' implies that the presence of 'a' triggers the '\_alt' replacement, regardless of other characters. This behavior resembles a keyboard shortcut encoder where lowercase 'a' is treated as a modifier key (Alt) when combined with other keys. I'll now hypothesize that the function replaces any occurrence of 'a' (case{-}insensitive) with '\_alt' and removes all other characters.",} \\
\texttt{  "hypothesis": "def f(text: str) {-}> str:\textbackslash{}n    \textbackslash{}"\textbackslash{}"\textbackslash{}"Replaces any occurrence of 'a' or 'A' with '\_alt' and removes all other characters.\textbackslash{}"\textbackslash{}"\textbackslash{}"\textbackslash{}n    return '\_alt' if 'a' in text.lower() else ''",} \\
\texttt{  "confidence\_score": 75} \\
\texttt{\}} \\
\texttt{```} \\
            }
        }
    }
    & & \\ \\

    \theutterance \stepcounter{utterance}  
    & & & \multicolumn{4}{p{0.6\linewidth}}{
        \cellcolor[rgb]{0.9,0.9,0.9}{
            \makecell[{{p{\linewidth}}}]{
                \texttt{\tiny{[P1$\langle$GM]}}
                \texttt{OUTPUT: (('aba', ''), True)} \\
            }
        }
    }
     \\ \\

    \theutterance \stepcounter{utterance}  
    & \multicolumn{4}{p{0.6\linewidth}}{
        \cellcolor[rgb]{0.561,0.737,0.561}{
            \makecell[{{p{\linewidth}}}]{
                \texttt{\tiny{[P1$\rangle$GM]}}
                \texttt{```json} \\
\texttt{\{} \\
\texttt{  "current\_mode": "probing",} \\
\texttt{  "input\_rationale": "The input 'aba' produces an empty string, which contradicts my previous hypothesis (which would have returned '\_alt'). This suggests the function's behavior is more nuanced than simply replacing any 'a' with '\_alt'. The pattern now appears to be: (1) 'a' alone → '', (2) 'aA' → '\_alt', (3) 'aba' → '', (4) 'a{-}b' → '\_alt'. This resembles a keyboard shortcut encoder where '\_alt' is only produced when there's exactly one lowercase 'a' combined with exactly one uppercase 'A' or certain other characters. I'll now hypothesize that the function returns '\_alt' only when there's exactly one 'a' (case{-}insensitive) combined with exactly one non{-}lowercase character, otherwise returns an empty string.",} \\
\texttt{  "hypothesis": "def f(text: str) {-}> str:\textbackslash{}n    \textbackslash{}"\textbackslash{}"\textbackslash{}"Returns '\_alt' if the text contains exactly one 'a' (case{-}insensitive) and exactly one non{-}lowercase character, otherwise returns empty string.\textbackslash{}"\textbackslash{}"\textbackslash{}"\textbackslash{}n    lower\_a\_count = text.lower().count('a')\textbackslash{}n    non\_lower\_count = sum(1 for c in text if not c.islower())\textbackslash{}n    return '\_alt' if lower\_a\_count == 1 and non\_lower\_count == 1 else ''} \\
            }
        }
    }
    & & \\ \\

    \theutterance \stepcounter{utterance}  
    & & & \multicolumn{4}{p{0.6\linewidth}}{
        \cellcolor[rgb]{0.9,0.9,0.9}{
            \makecell[{{p{\linewidth}}}]{
                \texttt{\tiny{[P1$\langle$GM]}}
                \texttt{Game Violation: ParseError: Response is not valid JSON.} \\
\texttt{Output ONLY:} \\
\texttt{One JSON with confidence\_score, current\_mode, input\_rationale, and hypothesis keys ...} \\
\texttt{OR} \\
\texttt{SOLVE: ```python ... ```} \\
            }
        }
    }
     \\ \\

    \theutterance \stepcounter{utterance}  
    & \multicolumn{4}{p{0.6\linewidth}}{
        \cellcolor[rgb]{0.561,0.737,0.561}{
            \makecell[{{p{\linewidth}}}]{
                \texttt{\tiny{[P1$\rangle$GM]}}
                \texttt{```json} \\
\texttt{\{} \\
\texttt{  "current\_mode": "confirming",} \\
\texttt{  "input\_rationale": "Based on all previous examples, I now see a clear pattern: the function returns '\_alt' when there is exactly one lowercase 'a' combined with exactly one uppercase letter (regardless of which one), and returns empty string in all other cases (including when there are multiple 'a's or no uppercase letters). The whitespace example shows it preserves other characters, but the key pattern is the 1:1 ratio of 'a' to uppercase letters. I'll make my final hypothesis with this logic.",} \\
\texttt{  "hypothesis": "def f(text: str) {-}> str:\textbackslash{}n    \textbackslash{}"\textbackslash{}"\textbackslash{}"Returns '\_alt' if text contains exactly one 'a' and exactly one uppercase letter, otherwise returns empty string.\textbackslash{}"\textbackslash{}"\textbackslash{}"\textbackslash{}n    a\_count = text.count('a')\textbackslash{}n    upper\_count = sum(1 for c in text if c.isupper())\textbackslash{}n    return '\_alt' if a\_count == 1 and upper\_count == 1 else ''",} \\
\texttt{  "confidence\_score": 95} \\
\texttt{\}} \\
\texttt{```} \\
            }
        }
    }
    & & \\ \\

    \theutterance \stepcounter{utterance}  
    & & & \multicolumn{4}{p{0.6\linewidth}}{
        \cellcolor[rgb]{0.9,0.9,0.9}{
            \makecell[{{p{\linewidth}}}]{
                \texttt{\tiny{[P1$\langle$GM]}}
                \texttt{OUTPUT: (('hello', '\_alt'), False)} \\
            }
        }
    }
     \\ \\

    \theutterance \stepcounter{utterance}  
    & \multicolumn{4}{p{0.6\linewidth}}{
        \cellcolor[rgb]{0.561,0.737,0.561}{
            \makecell[{{p{\linewidth}}}]{
                \texttt{\tiny{[P1$\rangle$GM]}}
                \texttt{```json} \\
\texttt{\{} \\
\texttt{  "current\_mode": "confirming",} \\
\texttt{  "input\_rationale": "The new example ('hello' → '\_alt') contradicts my previous hypothesis, which would have returned an empty string (no uppercase letters). This suggests the function's behavior is more complex. Looking at all examples: (1) 'a' → '', (2) 'aA' → '\_alt', (3) '   ' → '   ', (4) 'a{-}b' → '\_alt', (5) 'aba' → '', (6) 'hello' → '\_alt'. The pattern appears to be: return '\_alt' if there is exactly one lowercase 'l' in the string, otherwise return empty string (with special handling for the 'aA' case). However, this doesn't fully explain all cases. Given the time constraints, I'll make my final SOLVE attempt with the most consistent pattern I can identify: return '\_alt' if there's exactly one 'a' (case{-}insensitive) in the string, otherwise return empty string (with whitespace preserved)."} \\
\texttt{\}} \\
\texttt{```} \\
\\ 
\texttt{SOLVE:} \\
\texttt{```python} \\
\texttt{def f(text: str) {-}> str:} \\
\texttt{    """Returns '\_alt' if the input contains exactly one 'a' (case{-}insensitive), otherwise returns empty string (preserving whitespace)."""} \\
\texttt{    a\_count = text.lower().count('a')} \\
\texttt{    if a\_count == 1:} \\
\texttt{        return '\_alt'} \\
\texttt{    return text if} \\
            }
        }
    }
    & & \\ \\

    \theutterance \stepcounter{utterance}  
    & & & \multicolumn{4}{p{0.6\linewidth}}{
        \cellcolor[rgb]{0.9,0.9,0.9}{
            \makecell[{{p{\linewidth}}}]{
                \texttt{\tiny{[P1$\langle$GM]}}
                \texttt{Game Violation: ParseError: Response is not valid JSON.} \\
\texttt{Output ONLY:} \\
\texttt{One JSON with confidence\_score, current\_mode, input\_rationale, and hypothesis keys ...} \\
\texttt{OR} \\
\texttt{SOLVE: ```python ... ```} \\
            }
        }
    }
     \\ \\

    \theutterance \stepcounter{utterance}  
    & \multicolumn{4}{p{0.6\linewidth}}{
        \cellcolor[rgb]{0.561,0.737,0.561}{
            \makecell[{{p{\linewidth}}}]{
                \texttt{\tiny{[P1$\rangle$GM]}}
                \texttt{SOLVE:} \\
\texttt{```python} \\
\texttt{def f(text: str) {-}> str:} \\
\texttt{    """Returns '\_alt' if the input contains exactly one lowercase 'a' and any uppercase letter, otherwise returns empty string (preserving whitespace only)."""} \\
\texttt{    a\_count = text.count('a')} \\
\texttt{    has\_upper = any(c.isupper() for c in text)} \\
\texttt{    if a\_count == 1 and has\_upper:} \\
\texttt{        return '\_alt'} \\
\texttt{    return '' if not text.isspace() else text} \\
\texttt{```} \\
            }
        }
    }
    & & \\ \\

    \theutterance \stepcounter{utterance}  
    & & & \multicolumn{2}{p{0.3\linewidth}}{
        \cellcolor[rgb]{0.9,0.9,0.9}{
            \makecell[{{p{\linewidth}}}]{
                \texttt{\tiny{[GM$|$GM]}}
                \texttt{Game Verdict: LOSS} \\
            }
        }
    }
    & & \\ \\

    \theutterance \stepcounter{utterance}  
    & & & \multicolumn{2}{p{0.3\linewidth}}{
        \cellcolor[rgb]{0.9,0.9,0.9}{
            \makecell[{{p{\linewidth}}}]{
                \texttt{\tiny{[GM$|$GM]}}
                \texttt{Incorrect.} \\
\texttt{Failed Cases:} \\
\texttt{Input: {[}'aA'{]} | Expected:  | Got: \_alt} \\
\texttt{Input: {[}'hello world'{]} | Expected:   | Got: } \\
\texttt{Input: {[}'lk\_g5aIv'{]} | Expected:  | Got: \_alt} \\
            }
        }
    }
    & & \\ \\

    \theutterance \stepcounter{utterance}  
    & & & \multicolumn{2}{p{0.3\linewidth}}{
        \cellcolor[rgb]{0.9,0.9,0.9}{
            \makecell[{{p{\linewidth}}}]{
                \texttt{\tiny{[GM$|$GM]}}
                \texttt{def keep\_whitespace\_only(text: str) {-}> str:} \\
\texttt{    return ''.join(ch for ch in text if ch.isspace())} \\
            }
        }
    }
    & & \\ \\

    \theutterance \stepcounter{utterance}  
    & & & \multicolumn{2}{p{0.3\linewidth}}{
        \cellcolor[rgb]{0.9,0.9,0.9}{
            \makecell[{{p{\linewidth}}}]{
                \texttt{\tiny{[GM$|$GM]}}
                \texttt{def f(text: str) {-}> str:} \\
\texttt{    """Returns '\_alt' if the input contains exactly one lowercase 'a' and any uppercase letter, otherwise returns empty string (preserving whitespace only)."""} \\
\texttt{    a\_count = text.count('a')} \\
\texttt{    has\_upper = any(c.isupper() for c in text)} \\
\texttt{    if a\_count == 1 and has\_upper:} \\
\texttt{        return '\_alt'} \\
\texttt{    return '' if not text.isspace() else text} \\
            }
        }
    }
    & & \\ \\

    \theutterance \stepcounter{utterance}  
    & & & \multicolumn{2}{p{0.3\linewidth}}{
        \cellcolor[rgb]{0.9,0.9,0.9}{
            \makecell[{{p{\linewidth}}}]{
                \texttt{\tiny{[GM$|$GM]}}
                \texttt{Mode: passive\_labeled\_pairs} \\
\texttt{Accuracy: 0.840} \\
\texttt{Information requests: 5/15} \\
\texttt{Revealed passive examples: 6} \\
\texttt{Efficiency: 0.733} \\
\texttt{Slack: 10} \\
\texttt{Observed{-}pair consistency: 1.000} \\
\texttt{Membership balance: 3 true / 3 false} \\
\texttt{Parse errors: 4} \\
\texttt{Runtime errors: 0} \\
            }
        }
    }
    & & \\ \\

\end{supertabular}
}

\end{document}
